\documentclass[italian, a4paper, 10pt, twocolumn]{../../style/lab_unige}
\usepackage[a4paper, margin=1.75cm, footskip=0.25in]{geometry}

\usepackage[utf8]{inputenc}
\usepackage[T1]{fontenc}

\usepackage[italian]{babel}

% \usepackage{biblatex}

\usepackage[bookmarksopen=true, citebordercolor={0 1 0}, linkbordercolor={1 0 0}, urlbordercolor={0 1 1}]{hyperref}
\usepackage[numbered]{bookmark}

\usepackage{graphicx}
\graphicspath{{../fig/}}
\usepackage{array}
\usepackage{tabulary}
\usepackage{booktabs}

% FOUNDAMENTAL
\usepackage{../../style/custom}

\usepackage{physics}

\usepackage{breqn}
\usepackage{cuted}
\usepackage{txfonts}

\usepackage{lipsum}

%% Define ref types
\newcommand{\reftab}[1]{Tabella {\ref{#1}}}%
\newcommand{\reffig}[1]{Figura {\ref{#1}}}%
\newcommand{\refeqn}[1]{({\ref{#1}})}%
%% PAPER ONLY custom Macros


%%
\setlength{\columnsep}{6mm}

\begin{document}
    \twocolumn[
    \begin{@twocolumnfalse}
        \title{
            {\raggedright \includegraphics[width=0.2\linewidth]{../../style/lab_mark.pdf}\\}
            %%TITLE_HERE%%
        }
        \author{
        Eugenio Dormicchi\textsuperscript{1},
        % Riccardo Pizzimbone\textsuperscript{1}, 
        Giovanni Oliveri\textsuperscript{1},
        Mattia Sotgia\textsuperscript{1}
        }

        \date{
        \textsuperscript{1}Gruppo C03, Esperienza di laboratorio n. %%NN%% \\
        %\textsuperscript{2}In presenza in laboratorio per la presa dati\\
            % Università degli Studi di Genova, Dipartimento di Fisica.\\
            Presa dati-- 
            %%DATE_HERE%%, 15:00– 18:00; Analisi dati-- 
            <end-date here>
        }
        \maketitle
        
        \begin{abstract}
            \textit{Obiettivo-- }
        
            \textit{Metodi-- }
        
            \textit{Risultati-- }
        
            \textit{Conclusione-- }
        
        
        \end{abstract}
        \vspace{2em}
    \end{@twocolumnfalse}
    ]

    %%%% CORPO DEL TESTO
    %%%% CORPO DEL TESTO

    \section{Obiettivo}
    \label{section:aim}

    \section{Strumentazione}
    \label{section:strument}

    \section{Metodi}
    \label{section:methods}

    % \section{Analisi dati}
    % \label{section:analysis}

    \section{Risultati}
    \label{section:results}

    \section{Conclusione}
    \label{section:conclusion}

    \subsection{Controlli}

    \subsection{Possibili errori sistematici}

    \appendix

    \setcounter{table}{0}
    \renewcommand{\thetable}{A\arabic{table}}

    \section{Dati estesi}
    

\end{document}
    
